% eksperymentalne tłumaczenie na włoski

%

\section{Wyrocznia w Delfach} % lang-pl

\section{L'oracolo di Delfi} % lang-it
\index{wyrocznia w Delfach}

Przytoczymy teraz trzy problemy geometryczne, zaprzątające głowy przez blisko dwa tysiąclecia. % lang-pl

Wszystkie trzy znajdziemy też u Evesa \cite[s. 156]{eves1_1972}. % lang-pl
Presenteremo ora tre problemi geometrici che hanno occupato le menti dei matematici per quasi due millenni. % lang-it
Tutti e tre sono trattati anche da Eves \cite[p. 156]{eves1_1972}. % lang-it


Według legendy wyrocznia na wyspie Delos poradzi, by dla powstrzymania panującej tam zarazy trzeba dwukrotnie powiększyć ołtarz ofiarny, nie zmieniając jego sześciennego kształtu. % lang-pl
Secondo la leggenda, l'oracolo dell'isola di Delo consigliò, per fermare una pestilenza che vi infuriava, di raddoppiare il volume dell'altare sacrificale senza modificarne la forma cubica. % lang-it
\index{podwojenie sześcianu}

\begin{problem}[problem delijski] % lang-pl
\begin{problem}[problema di Delo] % lang-it
    \label{problem_delijski}%
    Zbudować sześcian o objętości dwa razy większej niż sześcian dany. % lang-pl
    
    Costruire un cubo avente volume doppio rispetto a un cubo assegnato. % lang-it
    \index{problem!delijski}
\end{problem}

https://www.deltami.edu.pl/2017/03/analiza-starozytnych-i-cyprian-norwid/

O problemie następne pokolenia dowiedzą się między innymi z prac Teona ze Smyrny. % lang-pl

Le generazioni successive conosceranno il problema anche grazie agli scritti di Teone di Smirne. % lang-it
\index[persons]{Teon ze Smyrny}% % lang-pl
\index[persons]{Teone di Smirne}% % lang-it
% https://en.wikipedia.org/wiki/Theon_of_Smyrna  It is also one of the sources of our knowledge of the origins of the classical problem of Doubling the cube.[5]
Choć odkryte zostanie wiele konstrukcji przybliżonych albo takich, które wykorzystują poza cyrklem i~linijką dodatkowe narzędzia (Hipokrates razem z Menaichmosem i hiperbola z parabolą, których nie da się wykreślić; Diokles i cysoida; Nikomedes i konchoida, Hippiasz z Elidy i kwadratrysa), to mimo upływu setek lat nie zostanie zrobiony żaden postęp. % lang-pl

Sebbene vengano scoperte molte costruzioni approssimate o costruzioni che utilizzano strumenti aggiuntivi oltre a riga e compasso (Ippocrate e Menecmo con l'iperbole e la parabola; Diocle con la cissoide; Nicomede con la concoide; Ippia di Elide con la quadratrice), nonostante il trascorrere dei secoli non si farà alcun progresso decisivo. % lang-it
\index[persons]{Hipokrates z Chios}% % lang-pl
\index[persons]{Menaichmos}% % lang-pl
\index{hiperbola}% % lang-pl
\index{parabola}%
\index[persons]{Diokles}% % lang-pl
\index{cysoida}% % lang-pl
\index[persons]{Nikomedes}% % lang-pl
\index{konchoida}% % lang-pl
\index[persons]{Hippiasz z Elidy}% % lang-pl
\index[persons]{kwadratrysa}% % lang-pl
\index[persons]{Ippocrate di Chio}% % lang-it
\index[persons]{Menecmo}% % lang-it
\index{iperbole}% % lang-it
\index[persons]{Diocle}% % lang-it
\index{cissoide}% % lang-it
\index[persons]{Nicomede}% % lang-it
\index{concoide}% % lang-it
\index[persons]{Ippia di Elide}% % lang-it
\index[persons]{quadratrice}% % lang-it

\begin{problem}[trysekcja kąta] % lang-pl
\begin{problem}[trisezione dell'angolo] % lang-it
    \label{trysekcja_kata}%
    Podzielić dowolny kąt na trzy równe części. % lang-pl
    
    Dividere un angolo arbitrario in tre parti uguali. % lang-it
    \index{trysekcja!kąta}
\end{problem}
% TODO: % TODO: Coxeter, s. 28
% TODO: https://en.wikipedia.org/wiki/Square_trisection
% https://en.wikipedia.org/wiki/Angle_trisection

https://www.deltami.edu.pl/2020/11/trysekcja-kata-w-geometrii-kartezjusza/

https://www.deltami.edu.pl/2015/09/zabawy-w-kacie/

Tu też pojawią się próby użycia różnorakich przyrządów. % lang-pl

Archmiedes zaproponuje dodanie podziałki do linijki, co okaże się wystarczające do wykreślenia kątów. % lang-pl
Anche qui verranno tentati approcci mediante strumenti di vario genere. % lang-it

Archimede proporrà di aggiungere una graduazione alla riga, il che si rivelerà sufficiente per la trisezione degli angoli. % lang-it
\index[persons]{Archimedes} % lang-pl
\index[persons]{Archimede} % lang-it
W średniowieczu zostanie zauważone, że zagadnienie trysekcji kąta prowadzi do rozwiązania równania algebraicznego stopnia trzeciego postaci % lang-pl

Nel Medioevo si osserverà che il problema della trisezione conduce alla risoluzione di un'equazione algebrica di terzo grado della forma % lang-it
\begin{equation}
    x^3 - 3 x - 2 \cos 3\alpha = 0.
\end{equation}
W szczególności, jeśli $\alpha = \pi / 9$, to pierwiastki nie są konstruowalne. % lang-pl

Patrz też do Coxetera \cite[s. 44]{coxeter_1967}. % lang-pl
In particolare, se $\alpha = \pi / 9$, le radici non sono costruibili. % lang-it
Si veda anche Coxeter \cite[p. 44]{coxeter_1967}. % lang-it
\begin{proposition} % lang-it
    La soluzione del problema di Delo e la trisezione dell'angolo non sono realizzabili con riga e compasso. % lang-it
\end{proposition} % lang-it


\begin{proposition} % lang-pl
    Rozwiązanie problemu delijskiego oraz trysekcja kąta są niewykonalne za pomocą cyrkla i linijki. % lang-pl
\end{proposition} % lang-pl
La prima dimostrazione di questo fatto fu data da Pierre Wantzel \cite{wantzel_1837} nel 1837. % lang-it

Pierwszy dowód tego faktu poda Pierre Wantzel \cite{wantzel_1837} w 1837 roku. % lang-pl
Być może to będzie powodem, dla którego dopiero w 1899 roku Frank Morley \cite{morley_1907} odkryje: % lang-pl
Forse anche per questo motivo soltanto nel 1899 Frank Morley \cite{morley_1907} scoprirà il seguente risultato: % lang-it
\index[persons]{Morley, Frank}

\begin{theorem}[Morleya]
    Punkty przeciecięcia tych spośród trójsiecznych kątów trójkąta, które sąsiadują z~którymś z~boków trójkąta, są wierzchołkami trójkąta równobocznego. % lang-pl
    
    I punti d'intersezione delle trisettrici degli angoli di un triangolo adiacenti ai suoi lati sono i vertici di un triangolo equilatero. % lang-it
    \index{trójkąt!Morleya}%
    \index{twierdzenie!Morleya}%
\end{theorem}

W tym sformułowaniu nie ma mowy, że chodzi o kąty wewnętrzne -- ponieważ twierdzenie jest prawdziwe także dla kątów zewnętrznych, z prawie takim samym dowodem; napisze o tym Kordos w $\Delta_{13}^{11}$. % lang-pl

Trójkąt równoboczny, jaki powstaje, ma bok długości % lang-pl
Nell'enunciato non si specifica che si tratta degli angoli interni, poiché il teorema è vero anche per gli angoli esterni, con una dimostrazione quasi identica; ne scrive Kordos in $\Delta_{13}^{11}$. % lang-it

Il triangolo equilatero ottenuto ha lato di lunghezza % lang-it
\begin{equation}
    8 R \sin \frac \alpha 3 \sin \frac \beta 3 \sin \frac \gamma 3,
\end{equation}
% Coxeter s. 23-25
% https://en.wikipedia.org/wiki/Hofstadter_points
a samo twierdzenie opiszą Coxeter \cite[s. 39-41]{coxeter_1967}, Zetel \cite[s. 177-179]{zetel_2020}, Neugebauer, Bogdańska \cite[s. 87, 105]{neugebauer_2018} (z dopiskiem, że nazwie się je kiedyś ,,ostatnim twierdzeniem geometrii elementarnej'' (?)). % lang-pl
e il teorema stesso è descritto da Coxeter \cite[p. 39-41]{coxeter_1967}, Zetel \cite[p. 177-179]{zetel_2020}, Neugebauer e Bogdańska \cite[p. 87, 105]{neugebauer_2018} (con l'osservazione che un giorno verrà chiamato ``l'ultimo teorema della geometria elementare'' (?)). % lang-it


Jak pamiętamy, pola powierzchnie wszystkich figur płaskich (trójkąta, równoległoboku, itd.) zostaną wyznaczone w starożytności przez podział figury na mniejsze części i ułożenie z nich prostokąta. % lang-pl
Come ricordiamo, le aree di tutte le figure piane (triangolo, parallelogramma ecc.) furono determinate nell'antichità suddividendole in parti più piccole e ricomponendole in un rettangolo. % lang-it

Próbowano tego samego wobec koła, bez powodzenia. % lang-pl
Si tentò di fare lo stesso con il cerchio, senza successo. % lang-it
\begin{problem}[quadratura del cerchio] % lang-it
\index{kwadratura koła}% % lang-it
    Costruire un quadrato avente area uguale a quella di un cerchio assegnato. % lang-it
\end{problem} % lang-it


\begin{problem}[kwadratura koła] % lang-pl
\index{kwadratura koła}% % lang-pl
    Wykreślić kwadrat o polu równym polu danego koła. % lang-pl
\end{problem} % lang-pl
Āpastamba, vissuto cinque secoli prima di Cristo e autore di uno dei quattro Śrautasūtra (testi vedici dedicati alla geometria), descrive come costruire un altare sacrificale. % lang-it

Apastamba, żyjący pięć wieków przed Chrystusem autor jednej z czterech śrautasutr (wedyjskiego tekstu poświęconemu geometrii), opisuje jak zbudować ołtarz ofiarny. % lang-pl
Oprócz próby kwadratury koła poda też przybliżenie $\sqrt 2$ z dokładnością do pięciu miejsc po przecinku. % lang-pl
Oltre a tentare la quadratura del cerchio, fornisce anche un'approssimazione di $\sqrt 2$ corretta a cinque cifre decimali. % lang-it


Opowiadaliśmy już o nie do końca trafionych pomysłach Mikołaja z Kuzy na stronie \pageref{nierownosc_mikolaja_z_kuzy}. % lang-pl
Abbiamo già parlato delle idee non del tutto corrette di Niccolò Cusano a pagina \pageref{nierownosc_mikolaja_z_kuzy}. % lang-it
\index[persons]{Mikołaj z Kuzy}% % lang-pl
\index[persons]{Niccolò Cusano}% % lang-it
% TODO: Nicolaus Cusanus, Nicolaus Krebs, Mikołaj Kuzańczyk
Adam Adamandy Kochański herbu Lubicz, dworzanin króla Jana Sobieskiego i członek zakonu jezuitów, znajdzie przybliżenie % lang-pl

Adam Adamandy Kochański dello stemma Lubicz, cortigiano del re Giovanni Sobieski e membro dell'ordine dei Gesuiti, troverà l'approssimazione % lang-it
\index[persons]{Kochański, Adam Adamandy}%
\begin{equation}
    \pi \approx \sqrt{\frac{40}{3} - 2 \sqrt{3}} = 3.14153\,33387...
\end{equation}
w pracy \emph{Observationes Cyclometricae ad facilitandam Praxin accommodatae} (Obserwacje cyklometryczne przystosowane dla ułatwienia praktycznego użycia) z 1685 roku. % lang-pl

Ludzie nie przestaną szukać innych niezbyt czasochłonnych przepisów. % lang-pl
nell'opera \emph{Observationes Cyclometricae ad facilitandam Praxin accommodatae} del 1685. % lang-it

Jacob de Gelder w 1849 znajdzie dokładniejsze przybliżenie % lang-pl
I matematici continueranno a cercare altre costruzioni semplici e pratiche. % lang-it

Nel 1849 Jacob de Gelder troverà un'approssimazione più accurata % lang-it
\index[persons]{Gelder@de Gelder, Jacob}%
\begin{equation}
    \pi \approx 3 + \frac{4^2}{7^2 + 8^2} = \frac{355}{113} = 3.14159\,29203...
\end{equation}

Ramanujan w 1914 roku poda inny przepis oparty o przybliżenie zgadzające się na ośmiu miejscach po przecinku: % lang-pl

Nel 1914 Ramanujan proporrà un'altra costruzione basata su un'approssimazione corretta a otto cifre decimali: % lang-it
\index[persons]{Ramanujan, Srinivasa}%
\begin{equation}
    \pi \approx \sqrt[4]{9^2 + \frac{19^2}{22}} = 3.14159\,26525...
\end{equation}

Konstrukcję Ramanujana przytoczy Jarosław Górnicki w $\Delta_{18}^3$. % lang-pl
La costruzione di Ramanujan è presentata da Jarosław Górnicki in $\Delta_{18}^3$. % lang-it


Ferdinand Lindemann odkryje w 1882 roku, że liczba $\pi$ pojawiająca się we wzorze na pole koła jest przestępna, zatem rozszerzenie $\mathbb Q(\sqrt{\pi}) / \mathbb Q$ jest za dużego stopnia i kwadratura koła, podobnie jak poprzednie dwa problemy, również nie jest wykonalna standardowymi narzędziami. % lang-pl
Nel 1882 Ferdinand Lindemann scoprirà che il numero $\pi$, che compare nella formula dell'area del cerchio, è trascendente; di conseguenza l'estensione $\mathbb Q(\sqrt{\pi})/\mathbb Q$ ha grado troppo elevato e la quadratura del cerchio, come i due problemi precedenti, non è realizzabile con gli strumenti classici. % lang-it

(Jeśli liczba $x$ jest konstruowalna, to stopień rozszerzenia $\mathbb Q(x) / \mathbb Q$ jest potęgą $2$, patrz \cite[s. 242]{hartshorne2000}; ale to nie wystarcza: wielomian $x^4 - x^3 - 5x^2 + 1$ ma cztery pierwiastki rzeczywiste $\alpha$ i rozszerzenia $\mathbb Q(\alpha)/\mathbb Q$ są stopnia $4$, ale grupa Galois ma rząd $12$ lub $24$, co nie jest potęgą dwójki). % lang-pl
(Se un numero $x$ è costruibile, allora il grado dell'estensione $\mathbb Q(x)/\mathbb Q$ è una potenza di $2$, vedi \cite[p. 242]{hartshorne2000}; ma questa condizione non è sufficiente: il polinomio $x^4-x^3-5x^2+1$ ha quattro radici reali $\alpha$ e le estensioni $\mathbb Q(\alpha)/\mathbb Q$ hanno grado $4$, mentre il gruppo di Galois ha ordine $12$ oppure $24$, che non sono potenze di due). % lang-it
\index[persons]{Lindemann, Ferdinand}

Kwadratura koła jest równoważna rektyfikacji okręgu, czyli zadaniu, by wykreślić odcinek tej samej długości, co obwód danego okręgu. % lang-pl

To również nie jest wykonalne. % lang-pl
La quadratura del cerchio è equivalente alla rettificazione del cerchio, cioè al problema di costruire un segmento avente la stessa lunghezza della circonferenza assegnata. % lang-it

Anche questo problema è impossibile. % lang-it

%